
\exercise
Show that the subspaces of $\R{2}$ are precisely $\{0\}$, all lines in $\R{2}$ containing the origin, and $\RR{2}$.


\exercise
Show that the subspaces of $\R{3}$ are precisely $\{0\}$, all lines in $\R{3}$ containing the origin, all planes in $\R{3}$ containing the origin, and $\RR{3}$.


\exercise
\begin{SUBtheorem}
\item
Let $U = \br[\big]{p \in \P\1_4(\F{}): p(6) = 0}$. Find a basis of $U$.
\item
Extend the basis in (a) to a basis of $\P\1_4(\F{})$.
\item
Find a subspace $W$ of $\P\1_4(\F{})$ such that $\P\1_4(\F{}) = U \oplus W\3$.
\end{SUBtheorem}


\exercise
\begin{SUBtheorem}
\item
Let $U = \br[\big]{p \in \P\1_4(\R{}): p'' \! (6) = 0}$. Find a basis of $U$.
\item
Extend the basis in (a) to a basis of $\P\1_4(\R{})$.
\item
Find a subspace $W$ of $\P\1_4(\R{})$ such that $\P\1_4(\R{}) = U \oplus W\3$.
\end{SUBtheorem}


\exercise \label{2basis}
\begin{SUBtheorem}
\item
Let $U = \br[\big]{p \in \P\1_4(\F{}): p(2) = p(5)}$. Find a basis of $U$.
\item
Extend the basis in (a) to a basis of $\P\1_4(\F{})$.
\item
Find a subspace $W$ of $\P\1_4(\F{})$ such that $\P\1_4(\F{}) = U \oplus W\3$.
\end{SUBtheorem}


\exercise
\begin{SUBtheorem}
\item
Let $U = \br[\big]{p \in \P\1_4(\F{}): p(2) = p(5) = p(6)}$. Find a basis of~$U$.
\item
Extend the basis in (a) to a basis of $\P\1_4(\F{})$.
\item
Find a subspace $W$ of $\P\1_4(\F{})$ such that $\P\1_4(\F{}) = U \oplus W\3$.
\end{SUBtheorem}


\exercise
\begin{SUBtheorem}
\item
Let $U = \br[\Big]{p \in \P\1_4(\R{}): \int_{-1}^1 p = 0}$. Find a basis of $U$.
\item
Extend the basis in (a) to a basis of $\P\1_4(\R{})$.
\item
Find a subspace $W$ of $\P\1_4(\R{})$ such that $\P\1_4(\R{}) = U \oplus W\3$.
\end{SUBtheorem}


\exercise
Suppose $v_1, \dots, v_m$ is linearly independent in $V$ and $w \in V\1$. Prove that
\[
\dim \Span(v_1 + w, \dots, v_m + w) \ge m-1.
\]
\exercisedisplayend


\exercise
Suppose $m$ is a positive integer and $p_0, p_1, \dots, p_m \in \P(\F{})$ are such that each
$p_k$ has degree~$k$.  Prove that $p_0, p_1, \dots, p_m$ is a basis of $\P\1_m(\F{})$.


\exercise
Suppose $m$ is a positive integer. For $0 \le k \le m$, let
\[
p_k(x) = x^k (1-x)^{m-k}.
\]
Show that $p_0, \ldots, p_m$ is a basis of $\P\1_m(\F{})$.
\exercisecomment{The basis in this exercise leads to what are called \marbf{Bernstein polynomials}. You can do a web search to learn how Bernstein polynomials are used to approximate continuous functions on\/ $[0,1]$.\index{Bernstein polynomials}}


\exercise
Suppose $U$ and $W$ are both four-dimensional subspaces of $\CC{6}$.
Prove that there exist two vectors in $U \cap W$ such that neither of these vectors is a scalar multiple of the other.


\exercise
Suppose that $U$ and $W$ are subspaces of $\R{8}$ such that $\dim U = 3$,
$\dim W = 5$, and $U + W = \RR{8}$.  Prove that $\R8 = U \oplus W\3$. \label{7exer}


\exercise
Suppose $U$ and $W$ are both five-dimensional subspaces of~$\RR{9}$.
Prove that $U \cap W \ne \{0\}$.


\exercise
Suppose $V$ is a ten-dimensional vector space and $V\1_1, V\1_2, V\1_3$ are subspaces of $V$ with
$\dim V\1_1 = \dim V\1_2 = \dim V\1_3 = 7$. Prove that $V\1_1 \cap V\1_2 \cap V\1_3 \ne \br{0}$.


\exercise
Suppose $V$ is finite-dimensional and $V\1_1, V\1_2, V\1_3$ are subspaces of $V$ with
$\dim V\1_1 + \dim V\1_2 + \dim V\1_3 > 2 \dim V\1$. Prove that $V\1_1 \cap V\1_2 \cap V\1_3 \ne \br{0}$.


\exercise
Suppose $V$ is finite-dimensional and $U$ is a subspace of $V$ with $U \ne V\1$. Let $n = \dim V$ and $m = \dim U$. Prove that there exist $n - m$ subspaces of $V\1$, each of dimension $n-1$, whose intersection equals $U$.


\exercise
Suppose that $V\1_1, \dots, V\1_m$ are finite-dimensional subspaces of~$V\1$. Prove that $V\1_1 + \dots + V\1_m$ is finite-dimensional and
\[
\dim (V\1_1 + \dots + V\1_m) \le \dim V\1_1 + \dots + \dim V\1_m.
\]
%\exercisedisplayend
\exercisecomment{The inequality above is an equality if and only if\/ $V\1_1 + \dots + V\1_m$ is a direct sum, as will be shown in \ref{3directsumdim}.}


\exercise
Suppose $V$ is finite-dimensional, with $\dim V = n \ge 1$.  Prove that there
exist one-dimensional subspaces $V\1_1, \dots, V\1_n$ of~$V$ such that \label{8exer}
\[
V = V\1_1 \oplus \dots \oplus V\1_n.
\]
\exercisedisplayend


\begin{comment} % Exercise deleted because it is redundant with \ref{3directsumdim}.
\exercise
Suppose that $V\1_1, \dots, V\1_m$ are finite-dimensional subspaces of~$V$ such that
$V\1_1 + \dots + V\1_m$ is a direct sum. Prove that \label{387}
\[
\dim V\1_1 \oplus \dots \oplus V\1_m = \dim V\1_1 + \dots + \dim V\1_m.
\]
\exercisedisplayend
\exercisecomment{This exercise deepens the  analogy between direct sums of
subspaces and disjoint unions of  subsets.  Specifically, compare this exercise to
the following statement: if a set is written as a disjoint union of
finite subsets, then the number of elements in the set equals the sum of the numbers of
elements in the disjoint subsets.}


\end{comment}

\exercise
Explain why you might guess, motivated by analogy with the formula for the number of elements in the
union of three finite sets, that if $V\1_1, V\1_2, V\1_3$ are subspaces
of a finite-dimensional vector space, then
\begin{align*}
\dim(V\1_1 + V\1_2 &+ V\1_3)\\
=&\dim V\1_1 + \dim V\1_2 + \dim V\1_3 \\
&{} - \dim (V\1_1 \cap V\1_2) - \dim (V\1_1 \cap V\1_3) - \dim (V\1_2 \cap V\1_3) \\
&+ \dim (V\1_1 \cap V\1_2 \cap V\1_3).
\end{align*}
Then either prove the formula above or give a counterexample.


\exercise
Prove that if $V\1_1, V\1_2$, and $V\1_3$ are subspaces of a finite-dimensional vector space, then
\begin{align*}
\dim&(V\1_1 + V\1_2+ V\1_3) \\[6bp]
= & \dim V\1_1 + \dim V\1_2 + \dim V\1_3\\[6bp]
&- \frac{ \dim(V\1_1 \cap V\1_2) + \dim(V\1_1 \cap V\1_3) + \dim(V\1_2 \cap V\1_3) }{3}\\[6bp]
&- \frac{ \dim\bigl( (V\1_1\! +\! V\1_2)\! \cap\! V\1_3 \bigr)
+ \dim\bigl( (V\1_1 \!+\! V\1_3)\! \cap\! V\1_2 \bigr)
+ \dim\bigl( (V\1_2\1 +\! V\1_3)\! \cap\! V\1_1 \bigr) }{3}.
\end{align*}
\exercisecomment{The formula above may seem strange because the right side does not look like an integer.}


